% from ega1-10-fr.tex
\oldpage[I]{193}

biunivoque canonique entre
$\operatorname{Hom}_{\mathfrak{S}}(\mathfrak{X},\mathfrak{Y})$ et
l'ensemble des suites $(u_n)$ de $S_n$-morphismes
$u_n:X_n\to Y_n$ rendant commutatifs les diagrammes
(\hyperref[I.10.6.7.1-fr]{10.6.7.1}).
\end{corollary}

\begin{proof}
Pour tout $\mathfrak{S}$-morphisme
$u:\mathfrak{X}\to\mathfrak{Y}$, on a par définition $f=g\circ u$, donc
\[
  u^*(\mathscr{L})\mathscr{O}_{\mathfrak{X}}
  =u^*(g^*(\mathscr{J})\mathscr{O}_{\mathfrak{Y}})
    \mathscr{O}_{\mathfrak{X}}
  =f^*(\mathscr{J})\mathscr{O}_{\mathfrak{X}}
  \subset\mathscr{K}
\]
le corollaire découle donc de
(\hyperref[I.10.6.9-fr]{10.6.9}).
\end{proof}

On notera que, pour $m\leq n$, la donnée d'un morphisme
$f_n:X_n\to Y_n$ détermine un morphisme et un seul
$f_m:X_m\to Y_m$ rendant commutatif le diagramme
(\hyperref[I.10.6.7.1-fr]{10.6.7.1}), comme on le voit aussitôt en se
ramenant au cas affine~; on a ainsi défini une application
$\varphi_{mn}:\operatorname{Hom}_{S_n}(X_n,Y_n)
\to\operatorname{Hom}_{S_m}(X_m,Y_m)$ et les
$\operatorname{Hom}_{S_n}(X_n,Y_n)$ forment pour les $\varphi_{mn}$ un
\emph{système projectif d'ensembles}~;
(\hyperref[I.10.6.11-fr]{10.6.11}) s'énonce encore en disant qu'il existe
une bijection canonique
\[
  \operatorname{Hom}_{\mathfrak{S}}(\mathfrak{X},\mathfrak{Y})
  \xrightarrow{\sim}
  \varprojlim_n\operatorname{Hom}_{S_n}(X_n,Y_n).
\]

\subsection{Produit de préschémas formels.}
\label{subsection:I.10.7-fr}

\begin{env}[10.7.1]
\label{I.10.7.1-fr}
Soit $\mathfrak{S}$ un préschéma formel~; les $\mathfrak{S}$-préschémas
formels formant une catégorie, on peut définir la notion de
\emph{produit} de $\mathfrak{S}$-préschémas formels.
\end{env}

\begin{proposition}[10.7.2]
\label{I.10.7.2-fr}
Soient $\mathfrak{X}=\operatorname{Spf}(B)$,
$\mathfrak{Y}=\operatorname{Spf}(C)$ deux schémas formels affines
au-dessus d'un schéma formel affine
$\mathfrak{S}=\operatorname{Spf}(A)$. Soient
$\mathfrak{Z}=\operatorname{Spf}(B\widehat{\otimes}_A C)$,
$p_1$, $p_2$ les $\mathfrak{S}$-morphismes correspondant
(\hyperref[I.10.2.2-fr]{10.2.2}) aux $A$-homomorphismes canoniques
(continus) $\rho$ et $\sigma$ de $B$ et $C$ dans
$B\widehat{\otimes}_A C$~; alors $(\mathfrak{Z},p_1,p_2)$ est un produit
des $\mathfrak{S}$-schémas formels affines $\mathfrak{X}$ et
$\mathfrak{Y}$.
\end{proposition}

\begin{proof}
En vertu de (\hyperref[I.10.4.6-fr]{10.4.6}), tout revient à vérifier que
si, à tout $A$-homomorphisme continu
$\varphi:B\widehat{\otimes}_A C\to D$, où $D$ est un anneau admissible qui
est une $A$-algèbre topologique, on associe le couple
$(\varphi\circ\rho,\varphi\circ\sigma)$, on définit une bijection
\[
  \operatorname{Hom}_A(B\widehat{\otimes}_A C,D)
  \xrightarrow{\sim}
  \operatorname{Hom}_A(B,D)\times\operatorname{Hom}_A(C,D)
\]
ce qui n'est autre que la propriété universelle du produit tensoriel
complété (\hyperref[0.7.7.6-fr]{0, 7.7.6}).
\end{proof}

\begin{proposition}[10.7.3]
\label{I.10.7.3-fr}
Étant donnés deux $\mathfrak{S}$-préschémas formels
$\mathfrak{X}$, $\mathfrak{Y}$, le produit
$\mathfrak{X}\times_{\mathfrak{S}}\mathfrak{Y}$ existe.
\end{proposition}

\begin{proof}
La démonstration est identique à celle de
(\hyperref[I.3.2.6-fr]{3.2.6}), en y remplaçant les schémas affines
(resp. les ouverts affines) par les schémas formels affines (resp. les
ouverts formels affines), et la prop.
(\hyperref[I.3.2.2-fr]{3.2.2}) par
(\hyperref[I.10.7.2-fr]{10.7.2}).
\end{proof}

Toutes les propriétés formelles du produit de préschémas
(\hyperref[I.3.2.7-fr]{3.2.7} et
\hyperref[I.3.2.8-fr]{3.2.8},
\hyperref[I.3.3.1-fr]{3.3.1} à
\hyperref[I.3.3.12-fr]{3.3.12}) sont valables sans aucune modification
pour le produit de préschémas formels.

\begin{env}[10.7.4]
\label{I.10.7.4-fr}
Soient $\mathfrak{S}$, $\mathfrak{X}$, $\mathfrak{Y}$ trois préschémas
formels et soient $f:\mathfrak{X}\to\mathfrak{S}$,
$g:\mathfrak{Y}\to\mathfrak{S}$ deux morphismes. Supposons qu'il existe
dans $\mathfrak{S}$, $\mathfrak{X}$, $\mathfrak{Y}$ respectivement,
trois systèmes fondamentaux d'Idéaux de définition
$(\mathscr{J}_\lambda)$, $(\mathscr{K}_\lambda)$,
$(\mathscr{L}_\lambda)$ respectivement, ayant même ensemble d'indices
$I$, tels que
$f^*(\mathscr{J}_\lambda)\mathscr{O}_{\mathfrak{X}}
\subset\mathscr{K}_\lambda$ et
$g^*(\mathscr{J}_\lambda)\mathscr{O}_{\mathfrak{Y}}
\subset\mathscr{L}_\lambda$ pour tout $\lambda$. Posons
$S_\lambda=(\mathfrak{S},
\mathscr{O}_{\mathfrak{S}}/\mathscr{J}_\lambda)$,
$X_\lambda=(\mathfrak{X},
\mathscr{O}_{\mathfrak{X}}/\mathscr{K}_\lambda)$,
$Y_\lambda=(\mathfrak{Y},
\mathscr{O}_{\mathfrak{Y}}/\mathscr{L}_\lambda)$~; pour
$\mathscr{J}_\mu\subset\mathscr{J}_\lambda$,
$\mathscr{K}_\mu\subset\mathscr{K}_\lambda$,
$\mathscr{L}_\mu\subset\mathscr{L}_\lambda$, notons que
$S_\lambda$ (resp. $X_\lambda$, $Y_\lambda$) est un sous-préschéma fermé
de $S_\mu$ (resp. $X_\mu$, $Y_\mu$) ayant même

